\section{A Type II estimate}\label{sec:typeii}

We apply the shifted-correlation estimate to factorizations $mn=2u+1$,
where $u$ carries a weight supported on many geometric prime bands.
The coefficient of $m$ will satisfy the discrepancy hypothesis of
\Cref{thm:shift}; the coefficient of $n$ may be arbitrary within its
stated bound. The task is to retain that discrepancy while converting
multiplicative factorizations into additive correlations.

We split $u=eh$ with $e$ slightly smaller than the scale of $m$, using
a smooth allocation of the band primes. Cauchy's inequality in $(e,n)$
then has a small enough outside factor to control its diagonal. Off the
diagonal, the two factorizations give cross-products whose difference is
a small nonzero integer. These are the shifted endpoints. All group
factors and the compulsory rough integer are assigned to $h$, so both
endpoints retain the form required by \Cref{thm:shift}.

\subsection{The weight and the uniform statement}

Use the prime groups and the weight $W_\ell$ of \Cref{sec:shifts}. In
particular, the groups are disjoint, their harmonic masses belong to a fixed
interval $[v_-,v_+]\subset(0,\infty)$, and there are between fixed positive
multiples of $T$ small groups and big groups. For fixed
$0<a<b<c<d<0.47$, the small primes lie in
$[\exp(L^a),\exp(L^b)]$ and the big primes in
$[\exp(L^c),\exp(L^d)]$. The big groups have the interval structure required
there. The damping $q\in(0,1)$ and the integer $\ell\geq1$ are fixed.
Write $u_*$ for the factor of $u$ supported outside all these groups and
$u_P=u/u_*$. Thus
\[
 W_\ell(u)=\prod_g
       q^{\omega_g(u)-\ell}\frac{(\omega_g(u))_\ell}{V_g^\ell}.
\]
All constants describing these groups are among the fixed data below.

Fix $0<c_1<c_2$, $s'>0$, and $\lambda>1$, with $c_2s'<1/4$.
For $0\leq j\leq j_x$ put
\[
 s_j=s'L\lambda^{-j},\qquad
 \mathcal Q_j=\{p\text{ prime}:c_1s_j\leq\log p\leq c_2s_j\}.
\]
Assume that these intervals are pairwise disjoint and disjoint from the
prime groups above, and that
$c_-T\leq s_{j_x}\leq c_+T$ for fixed $0<c_-\leq c_+$.
Choose a fixed integer $r_0$ such that
\begin{equation}\label{eq:typeii-band-room}
       (r_0-1)c_1/\lambda\geq c_2+1.
\end{equation}
There are $r_0$ labelled slots in every band. Each normally takes a prime
in $\mathcal Q_j$, with repetitions allowed. In one specified slot of a
fixed band $j_{**}$ replace the prime by an integer in
\begin{equation}\label{eq:typeii-rough-slot}
\begin{gathered}
 I_{**}=[x^{\tau_*},x^{\eta_*}],\qquad \Pminus(p)>W,\\
 \tau_*=c_1s'\lambda^{-j_{**}},\qquad
 \eta_*=c_2s'\lambda^{-j_{**}}.
\end{gathered}
\end{equation}
The notation $p$ for this slot will never assert that it is prime.
For sufficiently large $x$, this slot is free of group primes, since
$\exp(L^d)<W$.

Define
\begin{equation}\label{eq:typeii-weight}
\begin{aligned}
 a_Q(v)&=\left(\prod_{j=0}^{j_x}\frac1{r_0!}\right)
   \#\{\text{admissible labelled slot lists with product }v\},\\
 A(u)&=a_Q(u_*)W_\ell(u).
\end{aligned}
\end{equation}
The exceptional integer can be fixed as a divisor of $v$. Thereafter the
prime multiplicities in each disjoint band determine the list up to at
most $r_0!$ permutations. Consequently
\begin{equation}\label{eq:typeii-weight-bounds}
 0\leq a_Q(v)\leq\tau(v),\qquad
 0\leq A(u)\leq L^{C_A'}\tau(u_*),\qquad
 \sum_v\frac{a_Q(v)}v\leq L^{C_A'}
\end{equation}
for a fixed exponent $C_A'$. Here the bound for $W_\ell$ follows by
maximizing $q^{t-\ell}(t)_\ell/V_g^\ell$ in each group. For the last
bound, factor the harmonic sum slot by slot: every prime slot has bounded
harmonic mass, the exceptional slot has mass $O(L)$, and there are $O(T)$
slots. The exponent in \Cref{eq:typeii-weight-bounds} depends only on
the fixed data.

\begin{theorem}[Type II]\label{thm:typeii}
Fix $D_*>0$, $b_*>0$, and $C>0$, and require
$\eta_*<b_*/4$ in \Cref{eq:typeii-rough-slot}. Let $U,V$ be dyadic
scales satisfying
\[
 UV\asymp x,\qquad x^{b_*}\leq U,V\leq x^{1-b_*}.
\]
Let $\Psi$ be a fixed smooth function with compact support in a fixed
compact subinterval of $(0,\infty)$. Suppose that
$\abs{\alpha_m},\abs{\beta_n}\leq L^C$, and that $\alpha_m=0$ unless
$\Pminus(m)>W$. There is a fixed exponent $B$, depending only on these
data and the fixed group and band constants, such that the discrepancy
hypothesis in \Cref{eq:discrepancy} on $\alpha$ implies
\begin{equation}\label{eq:typeii-bound}
 \sum_{\substack{m\asymp U,\ n\asymp V,\ u\geq1\\mn=2u+1}}
      A(u)\Psi(u/x)\alpha_m\beta_n
       \ll xL^{-D_*}.
\end{equation}
Here the sequence to which the discrepancy hypothesis is applied includes
the actual restriction $m\asymp U$ and any additional interval restriction
on $m$. Explicitly, for that restricted sequence one requires
\[
 \left|\sum_{m\in I}\frac{\alpha_m\chi(m)m^{it}}m\right|
       \leq L^{-B}
 \quad\left(I\subset[1,x^2],\quad
       \operatorname{mod}(\chi)\leq L^B,\quad |t|\leq L^B\right).
\]
There is no discrepancy hypothesis on $\beta$. The constants are uniform
over the sequences, scales, and interval restrictions satisfying these
conditions.
\end{theorem}

\subsection{Removing a large group-prime factor}

\begin{proof}[Proof of \Cref{thm:typeii}]
We first discard $u_P>x^{b_*/4}$ with an error smaller than every fixed
negative power of $L$ times $x$. Set $\theta=L^{-d}$. For one group write
$a_p=(p^{1-\theta}-1)^{-1}$. Since $p^\theta\leq e$ and the group primes
tend to infinity, $a_p\leq 2e/p$. Its tilted marked harmonic sum is
\begin{align*}
 &\sum_{\substack{v:\ \text{all prime factors of }v\text{ in }\mathcal P_g}}
 \frac{q^{\omega_g(v)-\ell}(\omega_g(v))_\ell}
      {V_g^\ell v^{1-\theta}}\\
 &\hspace{1cm}=
 \frac1{V_g^\ell}\prod_{p\in\mathcal P_g}(1+qa_p)
 \sum_{\substack{p_1,\ldots,p_\ell\in\mathcal P_g\\
                         \text{distinct, ordered}}}
       \prod_{i=1}^{\ell}\frac{a_{p_i}}{1+qa_{p_i}}
 \leq C_0.
\end{align*}
The last constant is fixed because $v_-\leq V_g\leq v_+$.
Taking the product over $O(T)$ groups gives
\begin{equation}\label{eq:typeii-P-tail}
 \sum_{\substack{v>x^{b_*/4}\\v_*=1}}\frac{W_\ell(v)}v
       \leq \exp\bigl(-\tfrac14 b_*L^{1-d}\bigr)L^{C_1}.
\end{equation}
On the support of the sum in \Cref{eq:typeii-bound}, $u\asymp x$.
The number of divisors $m$ of $2u+1$ all of whose prime factors exceed
$W$ is at most $2^{\log(2u+1)/\log W}=\exp(O(\sqrt L))$.
Using \Cref{eq:typeii-weight-bounds,eq:typeii-P-tail}, multiplying
harmonic mass by $O(x)$, and using $1-d>1/2$, the total discarded
contribution is at most
\[
 xL^{C_2}\exp\bigl(-\tfrac14 b_*L^{1-d}+O(\sqrt L)\bigr).
\]
Call the remaining slot tuples good. This estimate is independent of
the splitting precision introduced next.

\subsection{A smooth partition into two factors}

Put $E=UL^{-K_0}$, where the fixed integer $K_0$ will be chosen
after the splitting costs have been bounded. Our target is a weighted
partition into factorizations $u=eh$ with
\begin{equation}\label{eq:typeii-e-range}
                 EL^{-C_s}\leq e\leq E,
\end{equation}
where $C_s$ depends only on the band constants. The reason for placing
$e$ below $U$ is that Cauchy's inequality will contribute an outside
factor $EV\asymp xL^{-K_0}$. We will use this logarithmic saving to
control the diagonal; it is therefore essential that the splitting
costs be independent of $K_0$.

Allocate all of $u_P$ and the exceptional integer slot to $h$. This
preserves the group weight and the long rough factor at each eventual
shifted endpoint. Process the remaining prime slots in increasing order
of their band index, and in a fixed order within each band. Write
$v_i=\log p_i$, and let $s_{j(i)}$ be the scale of slot $i$. If
$\delta_i\in\{0,1\}$ records whether it is allocated to $e$, put
\[
 R_i=\log E-\sum_{t<i}\delta_tv_t.
\]
Fix $\gamma\in C^\infty(\R)$ with $0\leq\gamma\leq1$,
$\gamma(t)=0$ for $t\leq0$, and $\gamma(t)=1$ for $t\geq1$.
At slot $i$ take it into $e$ with probability
$\gamma((R_i-v_i)/s_{j(i)})$; otherwise put it into $h$.

We verify the range in \Cref{eq:typeii-e-range} by tracking the
unfilled logarithmic size. At every stage the induction target is
\[
 0\leq R_i\leq\sum_{t\geq i}v_t+(c_2+1)s_{j_x}.
\]
The terminal allowance accounts for a skip in the final band.
On a good tuple the factors assigned to $h$ before the procedure starts
have product at most $x^{b_*/2}$. The sum of the available prime logarithms is therefore
at least $(1-b_*/2)L+O(1)$, whereas
$0<\log E\leq(1-b_*)L-K_0T$ for large $x$. This proves the initial
bound, even without the terminal allowance. Taking a slot of positive
probability requires $R_i>v_i$ and subtracts $v_i$ from both the residual
and the available logarithm, so it preserves the bound and nonnegativity.
Skipping a slot of positive probability requires
\begin{equation}\label{eq:typeii-skip}
              R_i<v_i+s_{j(i)}\leq(c_2+1)s_{j(i)}.
\end{equation}
If a later band exists, its unprocessed prime slots have total logarithm
at least $(r_0-1)c_1s_{j(i)+1}\geq(c_2+1)s_{j(i)}$ by
\Cref{eq:typeii-band-room}. The subtraction of one slot covers the
possible exceptional integer in that later band. Hence, after such a skip,
the residual fits inside the remaining available logarithm. In the final
band, a skip leaves at most $(c_2+1)s_{j_x}$, and subsequent operations
cannot increase this residual. This proves the induction target.
After processing the last slot it gives
\[
 0\leq\log E-\log e\leq(c_2+1)s_{j_x}.
\]
Since $s_{j_x}\leq c_+T$, this proves
\Cref{eq:typeii-e-range} with $C_s=(c_2+1)c_+$. The geometric series
for the remaining band scales also bounds $R_i/s_{j(i)}$ by a fixed
constant. Both bounds are independent of $K_0$.

Choose a fixed smooth compactly supported function $\rho$ with values in
$[0,1]$, equal to one throughout the resulting possible range of
$(R_i-v_i)/s_{j(i)}$. Replace the two transition functions by
\[
              f_1=\rho\gamma,\qquad f_0=\rho(1-\gamma).
\]
This changes no branch on a good tuple. For an arbitrary tuple the sum
of the two transitions is $\rho\leq1$, so its total branch mass is at
most one. Insert the resulting partition, with
\Cref{eq:typeii-e-range}, into \Cref{eq:typeii-weight}. Keep every
original factorial normalization. The branches sum to one on good
tuples, and extending back to all tuples costs at most the error in
\Cref{eq:typeii-P-tail}. Every resulting $h$ contains the exceptional
rough integer.

Here are quantitative details of the smooth separation, including its
independence from $K_0$. Let $N=O(T)$ be the number of processed slots,
and use the Fourier convention
\[
             f(t)=(2\pi)^{-1}\int\widehat f(\xi)e^{i\xi t}\,d\xi.
\]
Choose a fixed $C_F\geq1$ bounding the Fourier $L^1$ norms of both
$f_0,f_1$, with the normalization $(2\pi)^{-1}$ included.
Summing the Fourier total variations over the $2^N$ patterns gives
at most $(2C_F)^N=L^{O(1)}$. Truncating each $|\xi_i|$ at $L$ has
total error bounded, for every fixed $M$, by
\[
       2^NN C_M C_F^{N-1}L^{-M}.
\]
Thus any prescribed power saving is available. In a fixed pattern the
product of Fourier exponentials is a scalar of modulus one times
$\prod_t p_t^{i\sigma_t}$, where
\begin{equation}\label{eq:typeii-slot-frequencies}
 \sigma_t=-\frac{\xi_t}{s_{j(t)}}
                 -\delta_t\sum_{i>t}\frac{\xi_i}{s_{j(i)}}.
\end{equation}
Indeed the scalar is
$\exp(i\log E\sum_i\xi_i/s_{j(i)})$. Since
$\sum_i1/s_{j(i)}\ll1/T$, the frequencies in
\Cref{eq:typeii-slot-frequencies} are $O(L/T)$.
All the bounds in this paragraph are independent of $K_0$.

Separate $\Psi(eh/x)$ by Fourier inversion in $\log(eh/x)$ as well.
Its Fourier $L^1$ norm is fixed, and the tail outside frequency $L$
has arbitrary power saving; its two factors are powers of $e$ and $h$.
These errors can all be summed before Cauchy's inequality. Indeed, the
absolute sum over the original factorizations is bounded by
\begin{equation}\label{eq:typeii-absolute-pre-cauchy}
 L^{O(1)}\sum_{u\asymp x}\tau(u)\tau(2u+1)
                         \ll xL^{O(1)}
\end{equation}
by Cauchy's inequality and \Cref{lem:divisor-moments}.
The same estimate applies to each uniform transition or cutoff error,
using the total Fourier variations of the other factors. Thus the
original sum is, up to $O(xL^{-D_*-1})$, an integral and a sum of total
variation at most $L^{C_{f}}$ of expressions
\begin{equation}\label{eq:typeii-separated-B}
 B_0=\sum_{\substack{EL^{-C_s}\leq e\leq E\\n\asymp V}}
 a_e(e)\beta_n
       \sum_{\substack{m\asymp U,\ h\geq1\\mn=2eh+1}}
                     \alpha_m a_h(h).
\end{equation}
The exponent $C_f$ and every coefficient exponent in the next
display are independent of $K_0$:
\begin{equation}\label{eq:typeii-separated-coefficients}
\begin{gathered}
 \abs{a_e(e)}\leq L^{C_{a}},\qquad
 a_h(h)=h^{i\varrho}W_\ell(h)H(h_*),\\
 H(t)=\sum_{pr=t}
     \ind_{p\in I_{**},\,\Pminus(p)>W}p^{i\sigma}c_r,
 \quad \abs{c_r}\leq L^{C_{a}}.
\end{gathered}
\end{equation}
In particular $|a_h(h)|\leq L^{O(1)}\tau(h_*)$.
To verify these statements, only group-free prime slots go into $e$,
so $(eh)_*=eh_*$ and $W_\ell(eh)=W_\ell(h)$. For a fixed pattern,
the number of ordered pure-prime lists at a fixed product is at most
$(r_0!)^{j_x+1}=L^{O(1)}$. This proves both the $a_e$ bound and the
bound for the residual coefficient $c_r$. Fixing the exceptional integer
as a divisor gives the bound for $a_h$. Individual slot twists have
modulus one; the common power of $h$ from the cutoff is $h^{i\varrho}$.
The frequencies $\varrho,\sigma$ have fixed power bounds in $L$.

\subsection{Cauchy's inequality and the diagonal}

By Cauchy's inequality in $(e,n)$, \Cref{eq:typeii-separated-B} satisfies
\begin{equation}\label{eq:typeii-cauchy}
 |B_0|^2\leq EVL^{C_3}\mathcal N,
 \qquad
 \mathcal N=\sum_{e,n}\eta_e(e/E)\eta_n(n/V)
 \left|\sum_{\substack{m\asymp U,\ h\geq1\\mn=2eh+1}}
                           \alpha_ma_h(h)\right|^2.
\end{equation}
Here $\eta_e,\eta_n$ are nonnegative smooth majorants, bounded by a fixed
constant, equal to one on the original ranges. They vanish unless
\[
 \tfrac12 EL^{-C_s}\leq e\leq2E,\qquad n\asymp V.
\]
Their derivatives of order $j$ in their displayed arguments are
$O_j(L^{C_4(j+1)})$, with $C_3,C_4$ independent of $K_0$.
For instance the lower edge of $\eta_e$ is obtained by rescaling a
fixed smooth cutoff by $L^{C_s}$.

On the diagonal $m=m'$ the two equations force $h=h'$. For fixed $m,n$,
the allowable $e$ divide $(mn-1)/2$, and
\[
 \sum_{e\mid(mn-1)/2}\tau\bigl((mn-1)/(2e)\bigr)^2
                            \leq\tau(mn-1)^3.
\]
Collecting by $v=mn$ and using \Cref{lem:divisor-moments}, we obtain
\begin{equation}\label{eq:typeii-diagonal}
 \mathcal N_{\rm diag}
 \ll L^{O(1)}\sum_{v\asymp x}\tau(v)\tau(v-1)^3
 \ll xL^{C_5}.
\end{equation}
The exponents $C_3,C_5$ are independent of $K_0$. Since
$EV\asymp xL^{-K_0}$, its contribution to $|B_0|^2$ is
$O(x^2L^{-K_0+C_3+C_5})$. We can therefore fix $K_0$ large enough
that its square root, even multiplied by $L^{C_{f}}$, is
$O(xL^{-D_*-1})$. All subsequent precisions may depend on this fixed
$K_0$.

\subsection{The exact determinant parametrization}

Consider an off-diagonal term of \Cref{eq:typeii-cauchy}, with
\begin{equation}\label{eq:typeii-two-equations}
       mn=2eh+1,\qquad m'n=2eh'+1,\qquad m\ne m'.
\end{equation}
The shared variables $e,n$ force the two $m$ values to be congruent
modulo $2e$: the first equation makes $n$ a unit modulo $2e$, and
subtracting the equations gives $2e\mid m'-m$. Write
$m'-m=2ek$, with $k\ne0$. The resulting identities are
\begin{equation}\label{eq:typeii-determinant}
 m'-m=2ek,\qquad h'-h=kn,\qquad
 z=m'h,\qquad z+k=mh'.
\end{equation}
The last identity follows from
$mh'-m'h=k(mn-2eh)=k$. The enlarged $e$ range gives
\begin{equation}\label{eq:typeii-k-range}
                  0<|k|\ll L^{K_0+C_s}.
\end{equation}

Conversely, fix such a $k$, positive integers $m,m',h,h'$ with
$\Pminus(m),\Pminus(m')>W$, and representations
$m'h=z$, $mh'=z+k$. Suppose
\[
 m'\equiv m\pmod{2|k|},\qquad
                 e=\frac{m'-m}{2k}>0.
\]
Then $e$ is an integer and
\begin{equation}\label{eq:typeii-converse-identity}
                   m(h'-h)=k(2eh+1).
\end{equation}
Every prime factor of $m$ exceeds $W$, whereas $|k|$ is bounded by
a fixed power of $L$. Hence $(m,k)=1$ for sufficiently large $x$.
\Cref{eq:typeii-converse-identity} implies $m\mid2eh+1$.
Set $n=(2eh+1)/m$, a positive integer. The same identity gives
$h'-h=kn$, and substituting $m'=m+2ek$ recovers the second equation
of \Cref{eq:typeii-two-equations}. This proves a bijection, with all
variables recovered by the displayed formulas. In particular, no
additional divisibility condition is being dropped. The coprimality
$(m,k)=1$ is necessary to this converse.

\Cref{fig:determinant} displays the cross-pairing that produces the
two shifted endpoints.
\begin{figure}[htbp]
\centering
\begin{tikzpicture}[>=Stealth,every node/.style={font=\small},
  factor/.style={circle,draw=blue!50!black,fill=blue!3,
                 minimum size=7mm,inner sep=2pt}]
 \node[factor] (m) at (0,.9) {$m$};
 \node[factor] (h) at (4,.9) {$h$};
 \node[factor] (mp) at (0,-.9) {$m'$};
 \node[factor] (hp) at (4,-.9) {$h'$};
 \draw[->,gray!80!black] (m)--node[left]{$+2ek$}(mp);
 \draw[->,gray!80!black] (h)--node[right]{$+kn$}(hp);
 \draw[thick,blue!60!black] (m)--
   node[pos=.26,above,sloped]{$mh'=z+k$}(hp);
 \draw[thick,dashed,blue!60!black] (mp)--
   node[pos=.26,below,sloped]{$m'h=z$}(h);
 \node[align=left,anchor=west] at (6,0)
   {$mn-2eh=1$\\[3pt]
    $m'n-2eh'=1$\\[5pt]
    $mh'-m'h=k\ne0$};
\end{tikzpicture}
\caption{Forward determinant identities for factorizations with common
$e,n$. Cross-paired factors give $z=m'h$ and $z+k=mh'$; the vertical
arrows record signed increments, so either sign of $k$ is allowed.
The converse reconstruction, including its congruence, positivity,
and coprimality conditions, is proved in the text.}
\label{fig:determinant}
\end{figure}

Both $m$ and $m'$ are units modulo $q_k=2|k|$, so the congruence is
imposed exactly by
\begin{equation}\label{eq:typeii-character-congruence}
 \ind_{m'\equiv m\ (q_k)}
       =\frac1{\varphi(q_k)}\sum_{\chi\bmod q_k}
                                  \chi(m)\overline{\chi(m')}.
\end{equation}
Thus $\mathcal N_{\rm off}$ is a sum over $k,z$ and these characters
of terms having coefficient
\[
 \alpha_m\overline{\alpha_{m'}}\chi(m)\overline{\chi(m')}
                    a_h(h)\overline{a_h(h')},
 \qquad m'h=z,\quad mh'=z+k,
\]
with the enlarged smooth $(e,n)$ cutoffs. Their arguments are now the
exact functions
\begin{equation}\label{eq:typeii-pulled-back}
       e=\frac{m'-m}{2k},\qquad
       n=\frac{(m'-m)z}{kmm'}+\frac1m.
\end{equation}
The cutoffs are extended smoothly by zero to nonpositive arguments,
so they impose the positivity and size conditions as well.

\subsection{Smooth separation and the shifted endpoints}

On the support just obtained, $m,m'\asymp U$, $n\asymp V$, and
$h\asymp x/e$. Hence $z=m'h$ lies in
$[xL^{-C_6},xL^{C_6}]$ for a fixed $C_6$ now allowed to depend on
$K_0$. Insert a smooth dyadic partition in $z$ with $O(T)$ terms,
and denote one scale by $Z=xL^{O(1)}$. Keep its smooth cutoff
$\psi_Z(z/Z)$ outside the ensuing Fourier expansion.

For completeness, the pulled-back cutoffs have the following uniform
regularity. Put
$t_1=\log(m/U)$, $t_2=\log(m'/U)$, $t_3=\log(z/Z)$.
Multiply by fixed smooth localizations in these three variables equal
to one on the ranges in use. From \Cref{eq:typeii-pulled-back},
every derivative of $e/E$ in the $t_i$ is
$O_j(U/(|k|E))$, and every derivative of $n/V$ is
\[
                 O_j\bigl(Z/(|k|UV)+(UV)^{-1}\bigr).
\]
Both are bounded by fixed powers of $L$; moreover, the normalized
derivatives of $\eta_e,\eta_n$ have the bounds already given.
Repeated use of the chain rule therefore bounds derivatives of the
localized product cutoff by $O_j(L^{C_7(j+1)})$. Its support in
$(t_1,t_2,t_3)$ is a fixed compact box. Fourier inversion, as in
\Cref{lem:fourier-separation}, consequently represents this cutoff
with total variation $L^{O(1)}$ as a superposition of
\begin{equation}\label{eq:typeii-three-powers}
                        m^{i\xi_1}(m')^{i\xi_2}z^{i\xi_3}.
\end{equation}
Here is an explicit tail estimate. Enlarge a fixed exponent $C_8$ so
that, with $P=L^{C_8}$, the localized cutoff $F$ satisfies
$\|\partial^\alpha F\|_\infty\ll_\alpha P^{|\alpha|}$; its
zeroth derivative is bounded independently of $L$. Integration by parts
then gives
\[
 |\widehat F(\xi)|\ll_j(1+|\xi|/P)^{-j},\qquad
 \int_{\R^3}|\widehat F(\xi)|\,d\xi\ll P^3,\qquad
 \int_{|\xi|>PL}|\widehat F(\xi)|\,d\xi\ll_j P^3L^{3-j}.
\]
Thus truncation at $PL$ has uniform error $O(L^{-A})$ for every fixed
$A$, by choosing $j>A+3C_8+3$. These later exponents may depend on
$K_0$.

There is sufficient absolute control to sum these errors. If the
coefficients in \Cref{eq:typeii-separated-coefficients} and the two
$m$ sequences are replaced by their absolute values, each resulting
endpoint convolution is at most $L^{O(1)}\tau(z_*)^2$. Its product
with $W_\ell(z)$ has squared sum $O(ZL^{O(1)})$ on $z\asymp Z$.
Cauchy's inequality, also at $z+k$, therefore bounds the absolute
correlation by $ZL^{O(1)}$. This uses only
\Cref{lem:divisor-moments}. It remains true after dropping the
congruence; hence it controls the errors just described, after summing
the polynomial number of $k$ values, dyads, and character terms.

We spell out the endpoint identification to track exactly which
sequence satisfies discrepancy. Use the notation of
\Cref{eq:typeii-separated-coefficients} and fix a character $\chi$
and frequencies in \Cref{eq:typeii-three-powers}. Since $m,m'$ are
free of group primes,
\[
       z_*=m'h_*,\qquad (z+k)_*=mh'_*,\qquad
       W_\ell(z)=W_\ell(h),\quad W_\ell(z+k)=W_\ell(h').
\]
Define the invariant endpoint functions
\begin{align}
 F_0(v)&=\sum_{mpr=v_*}
   \alpha_m\chi(m)m^{i(\varrho-\xi_2)}
   \ind_{p\in I_{**},\,\Pminus(p)>W}p^{-i\sigma}\overline{c_r},
          \label{eq:typeii-first-endpoint}\\
 G_0(v)&=\sum_{mpr=v_*}
   \alpha_m\chi(m)m^{i(\xi_1+\varrho)}
   \ind_{p\in I_{**},\,\Pminus(p)>W}p^{-i\sigma}\overline{c_r}.
          \label{eq:typeii-second-endpoint}
\end{align}
Direct substitution shows that the separated summand is
\[
 \psi_Z(z/Z)z^{i(\varrho+\xi_3)}(z+k)^{-i\varrho}
       \overline{F_0(z)}G_0(z+k)W_\ell(z)W_\ell(z+k),
\]
up to scalar phases of modulus one and the Fourier coefficient.
For example, at $z=m'h$ the factor $h^{i\varrho}$ becomes
$z^{i\varrho}(m')^{-i\varrho}$, which explains the first endpoint's
power and the conjugation of $c_r$. In particular, the sequence inside
$F_0$ is the original $\alpha$, with its actual interval restriction.

Both \Cref{eq:typeii-first-endpoint,eq:typeii-second-endpoint} have
exactly the form in \Cref{eq:endpoint-form}. The character modulus
$2|k|$, all power frequencies, and the coefficient bounds are fixed
powers of $L$. The interval $I_{**}$ lies in
$[x^\tau,x^\eta]$ for fixed $0<\tau<\eta<1/4$, by
\Cref{eq:typeii-rough-slot} and the hypothesis $\eta_*<b_*/4$.
The residual $c_r$ has a fixed power bound, and no roughness condition
on $r$ is needed. The original $\alpha$ is supported on $W$-rough
integers and satisfies \Cref{eq:discrepancy}; the explicit character
and power factors are precisely the factors permitted in
\Cref{eq:endpoint-form}. Thus the discrepancy assumption has not been
silently transferred to a new coefficient sequence.

Finally set $\widetilde\psi_Z(t)=t\psi_Z(t)$. The last unweighted
correlation equals $Z$ times
\[
 \sum_{z>0,\ z+k>0}\frac{\widetilde\psi_Z(z/Z)}z
 z^{i(\varrho+\xi_3)}(z+k)^{-i\varrho}
       \overline{F_0(z)}G_0(z+k)W_\ell(z)W_\ell(z+k).
\]
This is \Cref{eq:shift-bound}, including its harmonic normalization.
Given any fixed $D_0$, \Cref{thm:shift} bounds it by $L^{-D_0}$
once the discrepancy exponent $B$ is sufficiently large. Consequently
each separated unweighted correlation is $O(ZL^{-D_0})$.
Choose $D_0$ after $K_0$, the $k$ range, all separation costs, and the
desired saving in \Cref{eq:typeii-cauchy}. The total off-diagonal
contribution is then $O(xL^{-A})$ for any prescribed fixed $A$.
Choose $B$ after this application of \Cref{thm:shift}. Together with
\Cref{eq:typeii-diagonal}, this gives
$|B_0|\ll xL^{-D_*-C_{f}-1}$. Summing its pre-Cauchy expansion
and the previously estimated tails proves \Cref{eq:typeii-bound}.
\end{proof}

\begin{remark}[Order of the precisions]\label{rem:typeii-order}
The band constants, $j_{**}$, $b_*$, coefficient bound $C$, and target
$D_*$ are fixed first. The transition functions, their total variation,
and the coefficient and diagonal exponents are independent of $K_0$;
only then is $K_0$ chosen. The smooth separation after Cauchy's
inequality and the accuracy required from \Cref{thm:shift} are fixed
next, and the required discrepancy exponent $B$ is fixed last.
Accordingly, when $\alpha$ is the difference between a test and its
rough-integer proxy, the precision of the cells defining that proxy may
be chosen after $B$, provided the proxy's coefficient bound is already
uniform in that cell precision. This is the order used in
\Cref{sec:primes}.
\end{remark}
